% Options for packages loaded elsewhere
\PassOptionsToPackage{unicode}{hyperref}
\PassOptionsToPackage{hyphens}{url}
%
\documentclass[
  11pt,
]{article}
\usepackage{amsmath,amssymb}
\usepackage{iftex}
\ifPDFTeX
  \usepackage[T1]{fontenc}
  \usepackage[utf8]{inputenc}
  \usepackage{textcomp} % provide euro and other symbols
\else % if luatex or xetex
  \usepackage{unicode-math} % this also loads fontspec
  \defaultfontfeatures{Scale=MatchLowercase}
  \defaultfontfeatures[\rmfamily]{Ligatures=TeX,Scale=1}
\fi
\usepackage{lmodern}
\ifPDFTeX\else
  % xetex/luatex font selection
\fi
% Use upquote if available, for straight quotes in verbatim environments
\IfFileExists{upquote.sty}{\usepackage{upquote}}{}
\IfFileExists{microtype.sty}{% use microtype if available
  \usepackage[]{microtype}
  \UseMicrotypeSet[protrusion]{basicmath} % disable protrusion for tt fonts
}{}
\makeatletter
\@ifundefined{KOMAClassName}{% if non-KOMA class
  \IfFileExists{parskip.sty}{%
    \usepackage{parskip}
  }{% else
    \setlength{\parindent}{0pt}
    \setlength{\parskip}{6pt plus 2pt minus 1pt}}
}{% if KOMA class
  \KOMAoptions{parskip=half}}
\makeatother
\usepackage{xcolor}
\usepackage[margin=1in]{geometry}
\setlength{\emergencystretch}{3em} % prevent overfull lines
\providecommand{\tightlist}{%
  \setlength{\itemsep}{0pt}\setlength{\parskip}{0pt}}
\setcounter{secnumdepth}{-\maxdimen} % remove section numbering
\ifLuaTeX
  \usepackage{selnolig}  % disable illegal ligatures
\fi
\IfFileExists{bookmark.sty}{\usepackage{bookmark}}{\usepackage{hyperref}}
\IfFileExists{xurl.sty}{\usepackage{xurl}}{} % add URL line breaks if available
\urlstyle{same}
\hypersetup{
  pdftitle={A literal-window construction-quality bound for Catalan's constant},
  pdfauthor={Anonymous working draft},
  hidelinks,
  pdfcreator={LaTeX via pandoc}}

\title{A literal-window construction-quality bound for Catalan's
constant}
\author{Anonymous working draft}
\date{12 August 2026}

\begin{document}
\maketitle
\begin{abstract}
We combine the Apéry-like row of Zudilin at index \(3n\) with the
\((4,7)\)-row of Nesterenko. After explicit unconditional integer
normalizations, both first coordinates contain the square of every prime
in the literal windows \(4n<\ell<6n\) and \(2n<\ell<3n\). One congruence
on the second coordinates therefore permits division of a combined row
by a modulus of logarithmic mass \(6n+o(n)\). The main new ingredient is
a two-successive-minima selection theorem. It supplies a lower bound for
the actual height, even when the congruence lattice has an exponentially
short first minimum, by reducing the second vector modulo the first in
the linear-form coordinate. The resulting integer approximations to
Catalan's constant have construction quality at least
\(0.5819834736933570952\ldots\). This is a statement about the displayed
construction, not an irrationality result; the corresponding
irrationality threshold is \(1\).
\end{abstract}

\hypertarget{statement-and-scope}{%
\section{1. Statement and scope}\label{statement-and-scope}}

Put

\[
 G=\beta(2)=\sum_{k=0}^{\infty}\frac{(-1)^k}{(2k+1)^2}.
\]

It is not known whether \(G\) is irrational. We therefore avoid the
phrase ``irrationality measure'' for the exponent studied here.

An \textbf{admissible approximation sequence} is a sequence
\((q_n,p_n)\in\mathbb Z^2\), defined for all sufficiently large \(n\),
such that

\[
 q_n\ne0,\qquad q_nG-p_n\ne0,\qquad |q_n|\longrightarrow\infty.
\]

Its \textbf{construction quality} is

\[
 \delta(q,p)=
 \liminf_{n\to\infty}
 \frac{-\log|G-p_n/q_n|}{\log|q_n|}.                 \tag{1.1}
\]

This definition deliberately measures the displayed integer
construction. In particular it is not the classical irrationality
exponent of an already-known irrational number. If

\[
 \log|q_n|=An+o(n),\qquad
 \log|q_nG-p_n|=En+o(n),\qquad A>0,
\]

then \(\delta(q,p)=1-E/A\). When \(E>0\), upper bounds for the height
and linear-form value alone are insufficient: a lower height bound is
essential.

The rows, modulus, and output pairs used in the theorem are defined in
Sections 2--4.

\begin{quote}
\textbf{Theorem 1.1.} For every sufficiently large \(n\), there is an
explicitly defined (in classical exact-real arithmetic) coefficient
vector \(c_n\in\mathbb Z^2\) for which the divided Zudilin--Nesterenko
combination \((q_n,p_n)\) is an admissible integer pair and \[
\boxed{\displaystyle
\delta(q,p)\ge
0.5819834736933570952\ldots .}                     \tag{1.2}
\] More exactly, the right side is \[
1-\frac{E_{\rm eff}}{A_{\rm eff}},\qquad
A_{\rm eff}=25.707995623896753637\ldots,\quad
E_{\rm eff}=10.746367029007697991\ldots .           \tag{1.3}
\]
\end{quote}

The proof uses only the literal prime windows displayed below. It does
not use the larger fractional-part window mass that occurs in other
denominator arguments.

\hypertarget{the-two-source-rows}{%
\section{2. The two source rows}\label{the-two-source-rows}}

Throughout, \(D_m=\operatorname{lcm}(1,\ldots,m)\).

\hypertarget{zudilins-row-at-index-3n}{%
\subsection{\texorpdfstring{2.1. Zudilin's row at index
\(3n\)}{2.1. Zudilin's row at index 3n}}\label{zudilins-row-at-index-3n}}

Let \(\varphi(t)=20t^2-8t+1\). Define \(Q_m,P_m\in\mathbb Q\) as the two
solutions of

\[
\begin{aligned}
 &(2m+1)^2(2m+2)^2\varphi(m)z_{m+1}\\
 &=(3520m^6+5632m^5+2064m^4-384m^3-156m^2+16m+7)z_m\\
 &\quad +(2m-1)^2(2m)^2\varphi(m+1)z_{m-1}\qquad(m\ge1),
\end{aligned}                                        \tag{2.1}
\]

with

\[
 (Q_0,Q_1)=(1,7/4),\qquad (P_0,P_1)=(0,13/8).
                                                               \tag{2.2}
\]

Write \(\lambda_m^Z=Q_mG-P_m\). Zudilin's 2003 theorem gives

\[
 2^{4m+3}D_mQ_m\in\mathbb Z,\qquad
 2^{4m+3}D_{2m-1}^3P_m\in\mathbb Z,                  \tag{2.3}
\]

and the partial-fraction proof in his 2002 remarks gives the additional
clearings

\[
 2^{6m}Q_m\in\mathbb Z,\qquad
 2^{6m}D_{2m-1}^2P_m\in\mathbb Z.                   \tag{2.4}
\]

Define

\[
 e_m=\min\{6m,\;4m+3+\lfloor\log_2(2m-1)\rfloor\}.   \tag{2.5}
\]

If a rational number is cleared both by \(2^a\) and by \(2^bM\), its
reduced denominator is a power of two of exponent at most
\(\min(a,b+v_2(M))\). Applying this observation to \(D_{2m-1}^2P_m\),
and similarly to \(Q_m\), proves that

\[
 X_n=2^{e_{3n}}D_{6n-1}^2Q_{3n},\qquad
 Y_n=2^{e_{3n}}D_{6n-1}^2P_{3n}                    \tag{2.6}
\]

are integers. The subexponential overhead is completely explicit:

\[
 2^{e_m}=2^{4m}s_m,\qquad
 s_m=2^{\min\{2m,\;3+\lfloor\log_2(2m-1)\rfloor\}},
 \qquad \log s_m=O(\log m).                          \tag{2.7}
\]

Set

\[
 \Lambda_n^Z=X_nG-Y_n.
                                                               \tag{2.8}
\]

\hypertarget{nesterenkos-47-row}{%
\subsection{\texorpdfstring{2.2. Nesterenko's
\((4,7)\)-row}{2.2. Nesterenko's (4,7)-row}}\label{nesterenkos-47-row}}

For \(n\ge1\), define

\[
 \mathcal R_n(t)=
 \frac{(3n)!(3/2)_{3n-1}}{(4n)!}
 \frac{t(t-1)\cdots(t-4n+1)}
      {\prod_{r=0}^{3n}(t+r+1/2)^2}.                 \tag{2.9}
\]

Define \(A_{1,n,j},A_{2,n,j}\in\mathbb Q\) uniquely by

\[
 \mathcal R_n(t)=\sum_{j=0}^{3n}
 \frac{(t+1)\cdots(t+j)}{(t+1/2)_{j+1}}
 \left(A_{1,n,j}+\frac{A_{2,n,j}}{t+j+1/2}\right).
                                                               \tag{2.10}
\]

Empty products are \(1\). Direct extraction of the double-pole
coefficient gives

\[
 A_{2,n,j}=2^{-14n+2j+1}
 \frac{(8n+2j)!\,j!\,(6n)!}
 {(4n)!\,(4n+j)!\,(3n-j)!^2\,(2j)!^2}>0.            \tag{2.11}
\]

Set

\[
 B_n=\sum_{j=0}^{3n}A_{2,n,j},                       \tag{2.12}
\]

\[
 C_n=\sum_{j=1}^{3n}\sum_{\nu=0}^{j-1}
 \frac{\nu!}{(3/2)_\nu}
 \left(A_{1,n,j}+\frac{A_{2,n,j}}{\nu+1/2}\right),   \tag{2.13}
\]

and

\[
 J_n=\int_0^1\!\int_0^1
 \frac{x^{4n-1/2}(1-x)^{3n}y^{4n}(1-y)^{3n-1/2}}
 {(1-xy)^{4n+1}}\,dx\,dy.                            \tag{2.14}
\]

Nesterenko's partial-fraction theorem gives

\[
 J_n=4B_nG-C_n.                                      \tag{2.15}
\]

The unconditional denominator theorem in the same paper implies

\[
 4^{7n}D_{6n}^2B_n\in\mathbb Z,\qquad
 4^{7n}D_{6n}^2C_n\in\mathbb Z.                     \tag{2.16}
\]

For clarity, the second inclusion can be read termwise from the proof.
Theorem 3 there gives \(4^{7n}A_{2,n,j}\in\mathbb Z\) and
\(4^{7n}D_{6n}A_{1,n,j}\in\mathbb Z\), after multiplication by the
harmless power \(4^j\). For \(0\le\nu\le3n-1\),

\[
 D_{6n}\frac{\nu!}{(3/2)_\nu}\in\mathbb Z,\qquad
 \frac{2D_{6n}}{2\nu+1}\in\mathbb Z,
\]

which clears the two summands in (2.13).

The integer row representing (2.15) is therefore

\[
 V_n=4^{7n+1}D_{6n}^2B_n,\qquad
 U_n=4^{7n}D_{6n}^2C_n,                              \tag{2.17}
\]

and

\[
 \Lambda_n^N=V_nG-U_n=4^{7n}D_{6n}^2J_n.            \tag{2.18}
\]

The factor \(4\) in \(V_n\) is forced by (2.15). A self-contained check
of this factor and of the sign in (2.11) is given in Appendix A.

\hypertarget{rates-the-literal-divisor-and-row-independence}{%
\section{3. Rates, the literal divisor, and row
independence}\label{rates-the-literal-divisor-and-row-independence}}

Put

\[
 \phi=\frac{1+\sqrt5}{2},\qquad
 T=\frac{3303+437\sqrt{57}}{144},\qquad
 T^-=\frac{3303-437\sqrt{57}}{144}.                  \tag{3.1}
\]

Define

\[
\begin{aligned}
 A_1&=12\log2+12+15\log\phi,\\
 E_1&=12\log2+12-15\log\phi,\\
 A_2&=14\log2+12+2\log T,\\
 E_2&=14\log2+12+2\log T^-.
\end{aligned}                                        \tag{3.2}
\]

The raw two-sided limits in the cited papers are

\[
\begin{array}{ll}
 n^{-1}\log Q_{3n}\to15\log\phi,&
 n^{-1}\log|\lambda_{3n}^Z|\to-15\log\phi,\\
 n^{-1}\log B_n\to2\log T,&
 n^{-1}\log J_n\to2\log T^-.
\end{array}                                          \tag{3.3}
\]

Together with \(\log D_m=m+o(m)\) and (2.7), these give the
exact-normalization limits

\[
\begin{array}{ll}
 \log X_n=A_1n+o(n),&\log Y_n=A_1n+o(n),\\
 \log V_n=A_2n+o(n),&\log U_n=A_2n+o(n),\\
 \log|\Lambda_n^Z|=E_1n+o(n),&
 \log|\Lambda_n^N|=E_2n+o(n).
\end{array}                                          \tag{3.4}
\]

The two second-coordinate statements are two-sided because
\(P_m/Q_m\to G>0\) and \(C_n/(4B_nG)\to1\). Numerically,

\[
\begin{aligned}
 A_1&=27.535943542613395425\ldots,&
 E_1&=13.099588790825292001\ldots,\\
 A_2&=29.354773862079159628\ldots,&
 E_2&=14.393145267190103982\ldots.
\end{aligned}                                        \tag{3.5}
\]

Define the literal-window modulus

\[
 S_n=
 \prod_{\substack{4n<\ell<6n\\\ell\ {\rm prime}}}\ell^2
 \prod_{\substack{2n<\ell<3n\\\ell\ {\rm prime}}}\ell^2.        \tag{3.6}
\]

Every prime in the two disjoint intervals occurs to the first power in
both \(D_{6n-1}\) and \(D_{6n}\). Such a prime is odd for \(n\ge1\). The
pure power-of-two clearing \(2^{18n}Q_{3n}\in\mathbb Z\) from (2.4)
shows that \(Q_{3n}\) is integral at every odd prime; likewise Theorem 3
in Nesterenko's paper shows that each \(A_{2,n,j}\), and hence \(B_n\),
is integral at these window primes after multiplication by a power of
two. Thus the displayed squares of the two lcm factors survive in the
integer first coordinates, and

\[
 S_n\mid X_n,\qquad S_n\mid V_n.                    \tag{3.7}
\]

Writing \(\vartheta(x)=\sum_{\ell\le x}\log\ell\), the prime number
theorem gives

\[
 \log S_n=
 2(\vartheta(6n)-\vartheta(4n))
 +2(\vartheta(3n)-\vartheta(2n))
 =6n+o(n).                                           \tag{3.8}
\]

The row determinant is

\[
 \mathcal D_n=X_nU_n-V_nY_n
              =V_n\Lambda_n^Z-X_n\Lambda_n^N.       \tag{3.9}
\]

The two terms have rates \(A_2+E_1\) and \(A_1+E_2\), and

\[
 \gamma=(A_2+E_1)-(A_1+E_2)
 =0.52527384310095222099\ldots>0.                   \tag{3.10}
\]

The reverse triangle inequality and (3.4) therefore give

\[
 \log|\mathcal D_n|=(A_2+E_1)n+o(n),                \tag{3.11}
\]

so the two rows are independent for all sufficiently large \(n\).

\hypertarget{the-congruence-lattice}{%
\section{4. The congruence lattice}\label{the-congruence-lattice}}

Define

\[
 L_n=\{(c_1,c_2)\in\mathbb Z^2:
       c_1Y_n+c_2U_n\equiv0\pmod {S_n}\}.            \tag{4.1}
\]

Let

\[
 g_n=\gcd(S_n,Y_n,U_n),\qquad d_n=S_n/g_n.           \tag{4.2}
\]

The homomorphism

\[
 \mathbb Z^2\longrightarrow\mathbb Z/S_n\mathbb Z,\qquad
 (c_1,c_2)\longmapsto c_1Y_n+c_2U_n
\]

has image generated by \(g_n\bmod S_n\), and hence

\[
 [\mathbb Z^2:L_n]=d_n.                             \tag{4.3}
\]

For \(c=(c_1,c_2)\in L_n\), define

\[
 q_n(c)=\frac{c_1X_n+c_2V_n}{S_n},\qquad
 p_n(c)=\frac{c_1Y_n+c_2U_n}{S_n}.                  \tag{4.4}
\]

These are integers by (3.7) and the defining congruence. Moreover,

\[
 q_n(c)G-p_n(c)
 =\frac{c_1\Lambda_n^Z+c_2\Lambda_n^N}{S_n}.        \tag{4.5}
\]

No estimate for the extra common divisor \(g_n\) will be needed.

\hypertarget{an-abstract-two-minimum-selection-theorem}{%
\section{5. An abstract two-minimum selection
theorem}\label{an-abstract-two-minimum-selection-theorem}}

We state the new ingredient for an arbitrary real number \(\alpha\). Let
\(X_n,Y_n,V_n,U_n\in\mathbb Z\), put

\[
 \lambda_{1,n}=X_n\alpha-Y_n,\qquad
 \lambda_{2,n}=V_n\alpha-U_n,
\]

and define \(L_n,q_n(c),p_n(c)\) by (4.1) and (4.4), with these four
integer rows in place of the Catalan rows. In this section, the six rate
statements (3.4) mean the corresponding assertions for
\(X_n,Y_n,V_n,U_n,\lambda_{1,n},\lambda_{2,n}\), and

\[
 \mathcal D_n=X_nU_n-V_nY_n
              =V_n\lambda_{1,n}-X_n\lambda_{2,n}.
\]

Put

\[
 s_n=\frac{\log S_n}{n},\qquad
 \kappa_n=\frac{\log d_n}{n},\qquad
 x_n=\frac{\kappa_n+E_2-E_1}{2},                    \tag{5.1}
\]

and

\[
 H_n=\kappa_n-x_n+A_2-s_n,\qquad
 F_n=x_n+E_1-s_n.                                   \tag{5.2}
\]

Thus

\[
 H_n-F_n=A_2-E_2,\qquad
 H_n+F_n=\kappa_n+A_2+E_1-2s_n.                    \tag{5.3}
\]

\begin{quote}
\textbf{Theorem 5.1 (two-minimum selection).} Suppose:

\begin{enumerate}
\def\labelenumi{\arabic{enumi}.}
\tightlist
\item
  \(S_n\mid X_n,V_n\), the lattice and output rows are given by
  (4.1)--(4.5), and \(d_n=[\mathbb Z^2:L_n]\);
\item
  the six two-sided rate statements (3.4) hold;
\item
  \(s_n\to K\), \(0\le\kappa_n\le s_n\), and \(A_2+E_1>A_1+E_2\);
\item
  \(H_n\) and \(F_n\) are eventually bounded away from zero.
\end{enumerate}

Then, for every sufficiently large \(n\), one can select \(c_n\in L_n\)
such that, with \(q_n=q_n(c_n)\), \(p_n=p_n(c_n)\), and
\(\ell_n=q_n\alpha-p_n\), \[
q_n\ne0,\qquad \ell_n\ne0,\qquad
\log|q_n|\ge H_nn-o(n),                            \tag{5.4}
\] and \[
\frac{\log|\ell_n|}{\log|q_n|}
\le\frac{F_n}{H_n}+o(1).                          \tag{5.5}
\]
\end{quote}

\textbf{Proof.} Consider the centrally symmetric rectangle

\[
 \mathcal C_n=\{(u,v)\in\mathbb R^2:
 |u|\le e^{x_nn},\ |v|\le e^{(\kappa_n-x_n)n}\}.    \tag{5.6}
\]

Its area is \(4d_n\). Let \(\mu_{1,n}\le\mu_{2,n}\) be the successive
minima of \(L_n\) relative to \(\mathcal C_n\). Minkowski's second
theorem in dimension two gives

\[
 \mu_{1,n}\mu_{2,n}\le1.                            \tag{5.7}
\]

Choose independent lattice vectors \(b_n^{(1)},b_n^{(2)}\) attaining the
successive minima. Their determinant is a nonzero multiple of \(d_n\),
while the rectangle bound and (5.7) give

\[
 |\det(b_n^{(1)},b_n^{(2)})|\le2d_n.
\]

Consequently

\[
 \log|\det(b_n^{(1)},b_n^{(2)})|=\log d_n+O(1).
                                                               \tag{5.8}
\]

Write

\[
 q_i=q_n(b_n^{(i)}),\quad p_i=p_n(b_n^{(i)}),\quad
 \ell_i=q_i\alpha-p_i.
\]

Choose \(\epsilon_n\downarrow0\) dominating all normalized remainders in
the rate statements and (5.8), and such that
\(n\sqrt{\epsilon_n}\to\infty\). Since \(\mu_{1,n}\le1\), put

\[
 r_n=-\frac{\log\mu_{1,n}}n\ge0.
\]

The definition of \(x_n\) balances the two linear-form contributions.
The second-row contribution to the first coordinate exceeds the
first-row contribution by the fixed crossed gap \((A_2+E_1)-(A_1+E_2)\).
Therefore

\[
\begin{array}{ll}
 |q_1|\le e^{(H_n-r_n+O(\epsilon_n))n},&
 |\ell_1|\le e^{(F_n-r_n+O(\epsilon_n))n},\\
 |q_2|\le e^{(H_n+r_n+O(\epsilon_n))n},&
 |\ell_2|\le e^{(F_n+r_n+O(\epsilon_n))n}.
\end{array}                                          \tag{5.9}
\]

The crossed-rate inequality and the two-sided rate assumptions give

\[
 \log|\mathcal D_n|=(A_2+E_1)n+o(n)
\]

by the reverse triangle inequality. Thus \(\mathcal D_n\ne0\), and in
particular the two source forms cannot both vanish, for all sufficiently
large \(n\). The determinant of the output rows is

\[
\begin{aligned}
 \Delta_n'&=q_1p_2-q_2p_1=q_2\ell_1-q_1\ell_2\\
 &=\frac{\det(b_n^{(1)},b_n^{(2)})\mathcal D_n}{S_n^2}.
\end{aligned}                                        \tag{5.10}
\]

Using this crossed-rate estimate, (5.3), and (5.8), we obtain

\[
 \log|\Delta_n'|=(H_n+F_n)n+O(\epsilon_nn).          \tag{5.11}
\]

There are three cases.

If \(r_n\le\sqrt{\epsilon_n}\), equations (5.9)--(5.11) force

\[
 \max(|q_1|,|q_2|)
 \ge e^{(H_n-O(\sqrt{\epsilon_n}))n}.
\]

Choose the vector with the larger first coordinate. If its linear form
vanishes, add or subtract the other vector with the sign for which the
first coordinates do not cancel. The two independent coefficient vectors
cannot both have zero linear form, since the two source forms are not
both zero. This gives

\[
 \log|q_n|=H_nn+o(n),\qquad
 \log|\ell_n|\le F_nn+o(n),                          \tag{5.12}
\]

and hence (5.4)--(5.5).

Suppose \(r_n>\sqrt{\epsilon_n}\) and \(\ell_1\ne0\). Choose
\(m_n\in\mathbb Z\) so that

\[
 0<|\ell_2-m_n\ell_1|\le\frac32|\ell_1|.             \tag{5.13}
\]

A nearest integer to \(\ell_2/\ell_1\) works unless the remainder is
zero; then use an adjacent integer. Set

\[
 c_n=b_n^{(2)}-m_nb_n^{(1)},\qquad
 q'=q_2-m_nq_1,\qquad \ell'=\ell_2-m_n\ell_1.
\]

Equation (5.10) gives the exact identity

\[
 q'=\frac{\Delta_n'}{\ell_1}
     +q_1\frac{\ell'}{\ell_1}.                       \tag{5.14}
\]

The first term has size at least \(e^{(H_n+r_n-O(\epsilon_n))n}\); the
second has size at most \(e^{(H_n-r_n+O(\epsilon_n))n}\). The first
therefore dominates by a factor tending to infinity. If
\(a_n=n^{-1}\log|\ell_1|\), then

\[
 n^{-1}\log|q'|=H_n+F_n-a_n+O(\epsilon_n),
\]

whereas \(a_n\le F_n-r_n+O(\epsilon_n)\) and
\(\log|\ell'|\le a_nn+O(1)\). Since \(a/(H_n+F_n-a)\) is increasing in
\(a\),

\[
 \frac{\log|\ell'|}{\log|q'|}
 \le\frac{F_n-r_n+O(\epsilon_n)}
          {H_n+r_n-O(\epsilon_n)}
 \le\frac{F_n}{H_n}+o(1).                           \tag{5.15}
\]

At the same time, \(\log|q'|\ge(H_n+r_n-O(\epsilon_n))n\), proving
(5.4).

Finally suppose \(\ell_1=0\). Then (5.10) implies \(q_1\ne0\) and
\(\ell_2\ne0\). Since \(x_n\), \(\kappa_n-x_n\), \(H_n\), and \(F_n\)
are bounded, a nonzero integral coordinate gives \(r_n=O(1)\). Thus
there are fixed constants \(B_q,B_\ell\) such that eventually
\(\log|q_2|\le B_qn\) and \(\log|\ell_2|\le B_\ell n\). Assumption 4 and
boundedness give \(\inf_nF_n/H_n>0\) after discarding finitely many
terms. Choose a fixed constant \(C>B_q\) so large that

\[
 \frac{B_\ell}{C}<\inf_n\frac{F_n}{H_n}.
\]

Then

\[
 c_n=b_n^{(2)}+\lceil e^{Cn}\rceil b_n^{(1)}
\]

has nonzero form \(\ell_2\), while \(\log|q_n(c_n)|\ge Cn-o(n)\). This
proves (5.4)--(5.5) in the last case. The case is included because
\(\ell_1=0\) would make \(\alpha\) rational, and no irrationality
assumption is permitted. \(\square\)

The theorem is canonical in classical exact-real arithmetic after fixing
a total order on \(\mathbb Z^2\) to resolve all choices of
successive-minimum vectors and ties.

\hypertarget{application-to-the-catalan-rows}{%
\section{6. Application to the Catalan
rows}\label{application-to-the-catalan-rows}}

For the present lattice,

\[
 s_n\to K=6,\qquad 0\le\kappa_n\le s_n.
\]

The quantities in (5.2) are uniformly positive. Moreover,

\[
\begin{aligned}
 \limsup_{n\to\infty}H_n
 &\le A_2-\frac{K+E_2-E_1}{2}=:A_{\rm eff},\\
 \limsup_{n\to\infty}F_n
 &\le E_2-\frac{K+E_2-E_1}{2}=:E_{\rm eff}.
\end{aligned}                                        \tag{6.1}
\]

Put

\[
 x_*=\frac{K+E_2-E_1}{2}
     =3.6467782381824059906\ldots.                   \tag{6.2}
\]

Then

\[
\begin{aligned}
 A_{\rm eff}&=A_2-x_*
 =25.707995623896753637\ldots,\\
 E_{\rm eff}&=E_2-x_*
 =10.746367029007697991\ldots.
\end{aligned}                                        \tag{6.3}
\]

Theorem 5.1 gives an admissible sequence and

\[
\begin{aligned}
 \delta(q,p)
 &=1-\limsup_{n\to\infty}
       \frac{\log|q_nG-p_n|}{\log|q_n|}\\
 &\ge1-\limsup_{n\to\infty}\frac{F_n}{H_n}.
\end{aligned}                                        \tag{6.4}
\]

The final step uses the exact relation \(H_n-F_n=A_2-E_2\), not a
comparison of separate limsups:

\[
 \frac{F_n}{H_n}=1-\frac{A_2-E_2}{H_n}.
\]

This expression is increasing in \(H_n>0\). Hence (6.1) gives

\[
 \delta(q,p)\ge
 \frac{A_2-E_2}{A_{\rm eff}}
 =1-\frac{E_{\rm eff}}{A_{\rm eff}}
 =0.5819834736933570952\ldots,                       \tag{6.5}
\]

which proves Theorem 1.1.

\hypertarget{comparison-and-literature-position}{%
\section{7. Comparison and literature
position}\label{comparison-and-literature-position}}

All comparisons in this section use Definition (1.1).

The square-cleared Zudilin row by itself has quality

\[
 1-\frac{E_1}{A_1}
 =0.5242731097791163788\ldots.                       \tag{7.1}
\]

The conservative unconditional Nesterenko row (2.17) has quality

\[
 1-\frac{E_2}{A_2}
 =0.5096829791700988894\ldots.                       \tag{7.2}
\]

Nesterenko's sharper \(\Delta_n\Delta_{1,n}\)-normalization has a
conjectural second-coordinate integrality. Under that condition, his
paper proves the safe exponent \(11/20=0.55\); its displayed limiting
log-ratio corresponds to a limiting construction quality near
\(0.55444\). The bound (6.5) exceeds each of these cited one-row
benchmarks.

A literature refresh through 12 August 2026 located the later
Krattenthaler--Zudilin hypergeometric paper and Zudilin's work on
simultaneous even beta values. Those papers supply important context but
do not state a better effective one-constant construction-quality bound
of the form used here. A 2022 Oberwolfach report still describes
Nesterenko's sequence as the best effective sequence for this
comparison. Because bibliographic searches are not proofs of priority,
we do not use the phrase ``best known'' in the theorem.

The decimal inequalities can be certified by rational interval
arithmetic. For \(y\in[1,2)\), put \(z=(y-1)/(y+1)\) and use

\[
 \log y=2\sum_{j=0}^{N}\frac{z^{2j+1}}{2j+1}+R_N,\qquad
 0<R_N<
 \frac{2z^{2N+3}}{(2N+3)(1-z^2)}.                   \tag{7.3}
\]

For example one may use

\[
\begin{aligned}
 \frac{2236067977499789}{10^{15}}&<\sqrt5<
 \frac{2236067977499790}{10^{15}},\\
 \frac{7549834435270749}{10^{15}}&<\sqrt{57}<
 \frac{7549834435270750}{10^{15}}.
\end{aligned}
\]

These brackets and \(N=40\) in (7.3) give, with outward rounding,

\[
 0.58198347<
 1-\frac{E_{\rm eff}}{A_{\rm eff}}
 <0.58198348.                                        \tag{7.4}
\]

\hypertarget{appendix-a.-the-two-printed-nesterenko-issues}{%
\section{Appendix A. The two printed Nesterenko
issues}\label{appendix-a.-the-two-printed-nesterenko-issues}}

The definitions (2.9)--(2.10), rather than the two problematic displayed
formulas in the printed paper, fix both conventions.

At \(t=-j-\tfrac12\), the numerator

\[
 t(t-1)\cdots(t-4n+1)
\]

has \(4n\) negative factors and is positive. The normalization in (2.9)
and the remaining denominator factors are positive. In the basis (2.10),
the sign contributed by

\[
 \frac{(t+1)\cdots(t+j)}{(t+1/2)_{j+1}}
 \frac1{t+j+1/2}
\]

to the double-pole coefficient is also positive: the numerator and the
denominator remaining after removal of \(t+j+1/2\) each have sign
\((-1)^j\). Thus \(A_{2,n,j}>0\), and direct factorial simplification
gives (2.11). The factor \((-1)^{4n+1}\) printed in Nesterenko's
diagonal formula (4.5) has the opposite sign.

Next, (2.15) implies identically

\[
 4^{7n}D_{6n}^2J_n
 =\bigl(4^{7n+1}D_{6n}^2B_n\bigr)G
  -4^{7n}D_{6n}^2C_n.
\]

Therefore the first coordinate must be \(V_n\) in (2.17). The factor
\(4\) is absent from the printed equation (4.4), but the paper uses
\(4^{7n+1}\) in its later sharper first coordinate. These are
normalization and sign corrections only; the integrality inclusions and
absolute-value asymptotics used above are unchanged.

\hypertarget{references}{%
\section{References}\label{references}}

\begin{enumerate}
\def\labelenumi{\arabic{enumi}.}
\tightlist
\item
  J. W. S. Cassels, \emph{An Introduction to the Geometry of Numbers},
  Springer, 1959.
\item
  C. Krattenthaler and W. Zudilin, \emph{Hypergeometry inspired by
  irrationality questions}, Kyushu J. Math. \textbf{73} (2019),
  189--203, doi:10.2206/kyushujm.73.189.
\item
  Yu. V. Nesterenko, \emph{On Catalan's constant}, Proc. Steklov Inst.
  Math. \textbf{292} (2016), 153--170, doi:10.1134/S0081543816010107.
\item
  W. Zudilin, \emph{A few remarks on linear forms involving Catalan's
  constant}, Chebyshevskii Sbornik \textbf{3}:2 (2002), 60--70,
  arXiv:math/0210423.
\item
  W. Zudilin, \emph{An Apéry-like difference equation for Catalan's
  constant}, Electron. J. Combin. \textbf{10} (2003), Research Paper 14,
  arXiv:math/0201024.
\item
  W. Zudilin, \emph{Arithmetic of Catalan's constant and its relatives},
  Abh. Math. Semin. Univ. Hambg. \textbf{89} (2019), 45--53,
  doi:10.1007/s12188-019-00203-w.
\item
  \emph{Rational approximations to Catalan's constant}, in Oberwolfach
  Report 21/2022.
\end{enumerate}

\end{document}
